\chapter{Родитель mixed superconnection curvature}
\label{ch:v10-h15-r2-superconnection-mixed-curvature-parent}

Обычное отображение момента прямой суммы не породило mixed bilinear.
Проверим более жёсткую конструкцию: две нечётные стрелки $A_\Sigma$ и
$\Sigma$ входят в один off-diagonal блок суперк связности, а требуемый
член читается из поляризации её квадрата.

На концах стрелки нейтральный $\Delta$-фон действует операторами
\begin{equation}
 L=T=6T_R^3,
 \qquad R=-B=-3(B-L).
\end{equation}
Для grading $\Gamma=\operatorname{diag}(I_8,-I_8)$ и нечётного поля
$\Phi(v)=\bigl(\begin{smallmatrix}0&D_v\\D_v&0\end{smallmatrix}\bigr)$
имеем $\{\Gamma,\Phi\}=0$. Поляризация нуль-форменной кривизны равна
\begin{equation}
 \frac12\{\Phi(A_\Sigma),\Phi(\Sigma)\}.
\end{equation}
Коэффициент $1/2$ компенсирует два порядка произведения и не является
новой ручкой.

Градуированный след фоновой кривизны даёт
\begin{equation}
 \operatorname{Str}\!\left[
 \operatorname{diag}(L,R)\,
 \frac12\{\Phi(A_\Sigma),\Phi(\Sigma)\}
 \right]
 =A_\Sigma^T(L-R)\Sigma
 =\boxed{A_\Sigma^TQ\Sigma},
\end{equation}
поскольку $L-R=T+B=Q=6Y$. Таким образом, и Cartan-направление, и общий
коэффициент возникают из одного квадрата суперк связности.

Однако положительный Hilbert--Schmidt trace складывает endpoint-блоки и
даёт
\begin{equation}
 A_\Sigma^T(L+R)\Sigma=A_\Sigma^T(T-B)\Sigma,
\end{equation}
а не $Q$. Разность правильного и обычного операторов равна $2B$ и имеет
rank $4$. Метрика суперследа $\Gamma$ имеет inertia $(8_-,0_0,8_+)$,
тогда как положительная trace metric имеет $(0_-,0_0,16_+)$. Поэтому
точный cross-член уже имеет каноническое алгебраическое происхождение, но
его включение в положительный полный curvature action требует выведенного
селекторного quotient или относительного следа.

Остаётся и прежняя опорная граница: условная стрелка $A_\Sigma$ имеет rank
$8$, но её унаследованное вложение в текущую суперк связность имеет rank
$0$.

\begin{theorem}[Точный graded-curvature источник]
Поляризация квадрата одной Real-нечётной суперк связности с endpoint moment
maps $(T,-B)$ точно порождает $A_\Sigma^TQ\Sigma$ с каноническим
коэффициентом. Обычный положительный след порождает вместо него $T-B$;
правильная разность существует только в градуированном trace channel.
\end{theorem}

\begin{nogocriterion}[Суперслед нельзя объявлять положительной нормой]
Нельзя считать physical parent закрытым, подменяя полный положительный
след суперследом лишь потому, что последний даёт правильный знак endpoint
moment maps. Требуется вывести relative curvature component, junk quotient
или другой положительный функционал, сохраняющий именно graded mixed block.
Также должна быть унаследована сама стрелка $A_\Sigma$.
\end{nogocriterion}

\begin{protocol}\begin{enumerate}
\item oddness и Real self-adjointness: \textbf{точно};
\item graded endpoint operator: \textbf{$L-R=Q$};
\item graded cross rank: \textbf{$8$};
\item coefficient origin: \textbf{канонический $1/2$ anticommutator};
\item positive-trace operator: \textbf{$L+R=T-B$};
\item positive-trace defect: \textbf{$2B$, rank $4$};
\item supertrace inertia: \textbf{$(8_-,0_0,8_+)$};
\item inherited/required auxiliary-arrow rank: \textbf{$0/8$};
\item conditional/physical origin: \textbf{$4/6$ и $4/6$};
\item следующий гейт: \textbf{аудит trace-селекторов curvature}.
\end{enumerate}\end{protocol}

Машинный сертификат сохранён в
\path{s2t_v10_particle_wrinkle_dislocation_callias_equal_charge_twist_m4_h15_r2_superconnection_mixed_curvature_parent_origin_gate_results.json}.